\chapter{2008:Osamu Shimomura and Martin Chalfie and Roger Y. Tsien}

\label{chap:chemistry2008}

The Nobel Prize in Chemistry for the year 2008 was awarded jointly to Osamu Shimomura and Martin Chalfie and Roger Y. Tsien.

\textbf{Motivation:}

for the discovery and development of the green fluorescent protein, GFP

\section{Osamu Shimomura}

\subsection*{Early Life and Education}

Osamu Shimomura was born in 1928 in Fukuchiyama, Japan.

\subsection*{Career and Major Research}

Shimomura studied at Nagasaki Medical College (later part of Nagasaki University) and moved to the United States, where he worked at Princeton University and Boston University and carried out research at the Marine Biological Laboratory in Woods Hole. In the early 1960s he isolated the bioluminescent protein aequorin and the green fluorescent protein (GFP) from the jellyfish Aequorea victoria.

\subsection*{Nobel Prize-winning Work}

The work for which Osamu Shimomura was awarded the Nobel Prize in Chemistry in 2008 was described by the Nobel Committee as: "for the discovery and development of the green fluorescent protein, GFP".

Shimomura discovered GFP and showed that it glows green under ultraviolet light on its own, without any additional protein or cofactor.

\subsection*{Significance and Impact}

Because GFP's light-emitting group is self-contained, the protein made it possible to attach a visible tag to other proteins inside living cells. This discovery provided the tool that later developments turned into a universal marker for biology.

\subsection*{Later Career and Honors}

Shimomura was professor emeritus of Boston University and continued his work on bioluminescent proteins. He died on 2018-10-19 in Nagasaki, Japan.

\subsection*{Legacy}

The discovery of GFP opened the way for imaging gene expression and protein behavior in living organisms, a technique now used across biology and medicine.

\subsection*{Modern Developments}

Shimomura's discovery of the green fluorescent protein (GFP) and its mechanism of luminescence founded the field of in vivo imaging. GFP and its engineered variants are now used in thousands of laboratories to visualize proteins, cells, and gene expression in living organisms, transforming cell biology and neuroscience.

\subsection*{Selected Publications}

"Extraction, Purification and Properties of Aequorin, a Bioluminescent Protein from the Luminous Hydromedusan, Aequorea" (Journal of Cellular and Comparative Physiology, 1962)

\clearpage

\section{Martin Chalfie}

\subsection*{Early Life and Education}

Martin Chalfie was born in 1947 in Chicago, Illinois, USA.

\subsection*{Career and Major Research}

Chalfie earned bachelor's and doctoral degrees from Harvard University, in 1969 and 1977, and did postdoctoral work with Sydney Brenner at the MRC Laboratory of Molecular Biology in Cambridge, where he studied the nervous system of the nematode Caenorhabditis elegans. He became a professor at Columbia University.

\subsection*{Nobel Prize-winning Work}

The work for which Martin Chalfie was awarded the Nobel Prize in Chemistry in 2008 was described by the Nobel Committee as: "for the discovery and development of the green fluorescent protein, GFP".

Chalfie demonstrated that the gene for GFP could be inserted into other organisms and that the protein would light up wherever the gene was active, providing a direct readout of gene expression.

\subsection*{Significance and Impact}

His demonstration made GFP the marker of choice for following gene expression and protein location in living cells. It showed that GFP works as a genetic marker in a wide range of organisms.

\subsection*{Later Career and Honors}

Chalfie continues his research on the nematode nervous system at Columbia University.

\subsection*{Legacy}

Chalfie's work established the standard way of visualizing gene activity in living organisms, from bacteria to mammals.

\subsection*{Modern Developments}

Chalfie's demonstration that GFP could be used as a genetic marker to visualize gene expression in living organisms established GFP as the tool of choice for biological imaging. His work made possible the study of neurons, development, and disease in transparent organisms and founded modern optogenetics and live-cell imaging.

\subsection*{Selected Publications}

"Green Fluorescent Protein as a Marker for Gene Expression" (Science, 1994)

\clearpage

\section{Roger Y. Tsien}

\subsection*{Early Life and Education}

Roger Y. Tsien was born in 1952 in New York City, New York, USA.

\subsection*{Career and Major Research}

Tsien earned a bachelor's degree from Harvard University in 1972 and a doctorate from the University of Cambridge in 1977. He became professor of pharmacology, chemistry, and biochemistry at the University of California, San Diego, and an investigator of the Howard Hughes Medical Institute.

\subsection*{Nobel Prize-winning Work}

The work for which Roger Y. Tsien was awarded the Nobel Prize in Chemistry in 2008 was described by the Nobel Committee as: "for the discovery and development of the green fluorescent protein, GFP".

Tsien worked out how the GFP chromophore forms and engineered variants that glow in different colors, from blue to red, and designed fluorescent indicators that report the activity of cells.

\subsection*{Significance and Impact}

His colored variants allowed several different proteins to be tracked at the same time inside one cell. Tsien's palette of fluorescent proteins is now used throughout cell biology and neuroscience.

\subsection*{Later Career and Honors}

Tsien worked at the University of California, San Diego, for most of his career. He died on 2016-08-24 in Eugene, Oregon, USA.

\subsection*{Legacy}

The rainbow of fluorescent proteins he created underlies modern imaging of the interior of living cells, including studies of the brain.

\subsection*{Modern Developments}

Tsien's engineering of GFP variants with new colors and improved properties created a palette of fluorescent proteins that revolutionized biological imaging. His techniques underpin modern super-resolution microscopy, cell tracking, and the study of neuronal activity, and his fluorescent indicators remain essential tools in neuroscience and cell biology.

\subsection*{Selected Publications}

"The Green Fluorescent Protein" (Annual Review of Biochemistry, 1998)

"Improved Green Fluorescence" (Nature, 1995)

\clearpage

\section*{Chapter Summary}

The 2008 prize recognized the discovery and development of the green fluorescent protein: Shimomura found GFP in a jellyfish, Chalfie made it a marker of gene expression, and Tsien created color variants for imaging. GFP-based imaging is now among the most widely used tools in biology.
